\section{Introduction}
\label{sec:introduction}

Diabetes mellitus affects approximately 537 million adults worldwide, a
figure projected to reach 783 million by 2045
\citep{idf-atlas-10th-edition-2021}. In Israel, the national diabetes
registry documents a steadily growing patient population over the past
decade
\citep{israel-national-center-disease-control-diabetes-registry-report-2012-2022},
with a corresponding burden on the health system
\citep{summary-cost-diabetes-israel-2018}. That burden is concentrated in
hospital care: approximately 40\% of patients admitted to Israeli general
hospitals have diabetes \citep{hochberg2023inpatient}.

In hospital the clinically relevant object is not the diagnosis but the
glycemic state during the admission. Hyperglycemia arising in hospital is
associated with infection, delayed healing, metabolic decompensation and
increased mortality, and occurs in patients with and without documented
diabetes. Its counterpart, hypoglycemia, is acute, harmful, and largely
iatrogenic---a consequence of the insulin used to correct the first problem.
Inpatient glycemic management is therefore a balancing act shaped less by any
single decision than by how monitoring, medication management and insulin
dosing are carried out repeatedly across a stay. That makes the treatment
process itself, not only the patient's state, worth measuring.

Guidelines exist to standardize such processes. In September 2023 the Israel
Endocrine Society, the Israel Diabetes Association and the National Diabetes
Council issued a national position paper on inpatient diabetes care
\citep{hochberg2023inpatient}, specifying how often glucose should be
measured, how oral antidiabetic agents should be handled on admission, when
basal and correction insulin should be given, and how to respond to
hyperglycemic and hypoglycemic episodes. Recommendations of this kind are
checkable in principle because they are procedural, temporal and
conditional: an action is expected, within some window or at some frequency,
when a triggering situation holds.

Establishing whether they were followed is harder. Guidelines are published
as text, must be adapted to the individual patient, and are known to be
applied incompletely. Retrospective assessment has conventionally required
chart review, which does not scale past modest single-institution
samples---too small to separate one ward's practice from the effect of a
guideline, and far too small to see how practice shifts when a guideline is
introduced nationally. The question that matters most to a guideline's
authors is therefore rarely answerable: not whether the guideline is sound,
but which of its individual recommendations are adopted, and which of them
measurably change what happens to patients.

This study addresses that question. Each recommendation was formalized as a
temporal pattern over abstracted clinical data, using knowledge-based
temporal abstraction \citep{shahar1997framework} and the quality-assessment
formulation of adherence \citep{ben2024supporting}, producing a continuous
compliance score automatically wherever the recommendation applies. Applying
this to 48,354 admissions across six general hospitals spanning the
guideline's introduction, we ask how completely each recommendation was
implemented and whether that changed after issue; whether adherence is
associated with subsequent outcomes once measured patient characteristics
are accounted for; and which individual recommendations carry that
association. Adherence thus serves as an explanatory variable, and the
guideline's recommendations as hypotheses to be tested against what
subsequently happened to the patient.

% JBI requires a Statement of Significance at the end of the Introduction.
\begin{table}[h]
\centering
\caption{Statement of Significance}
\label{tab:significance}
\begin{tabularx}{\textwidth}{
>{\raggedright\arraybackslash}p{0.22\textwidth}
>{\raggedright\arraybackslash}X
}
\toprule

\textbf{Problem} &
Adherence to inpatient diabetes guidelines is known to be incomplete, but
assessing it retrospectively has required manual chart review. This does not
scale to the size needed to distinguish local practice from the effect of a
guideline, or to observe how practice changes when one is introduced
nationally. \\

\midrule

\textbf{What is Already Known} &
Knowledge-based temporal abstraction and quality-assessment temporal
patterns allow guideline adherence to be scored automatically from clinical
records, and have been demonstrated in single-institution settings. Prior
applications evaluate care \emph{against} the guideline, which serves as the
reference standard. \\

\midrule

\textbf{What This Paper Adds} &
An automated, component-level assessment of a national inpatient
hyperglycemia guideline across six hospitals and 48,354 admissions spanning
its introduction, together with a transferable knowledge base encoding it.
By treating each recommendation as a testable hypothesis rather than as a
fixed standard, the study identifies which recommendations were least
implemented and which are associated with subsequent outcomes---and finds
these to be largely the same set. \\

\bottomrule
\end{tabularx}
\end{table}
